\section{AMC 12 12A 2009}
        \begin{enumerate}
        	\item (\textit{AMC 12 12A 2009}) Kim's flight took off from Newark at 10:34 AM and landed in Miami at 1:18 PM. Both cities are in the same time zone. If her flight took $h$ hours and $m$ minutes, with $0 < m < 60$, what is $h + m$?


 $\textbf{(A)}\ 46 \qquad \textbf{(B)}\ 47 \qquad \textbf{(C)}\ 50 \qquad \textbf{(D)}\ 53 \qquad \textbf{(E)}\ 54$

 \textbf{Answer: }A



\textbf{Solution}

There is $1$ hour and $60-34 = 26$ minutes between 10:34 AM and noon; and there is $1$ hour and $18$ minutes between noon and 1:18 PM. Hence the flight took $2$ hours and $26 + 18 = 44$ minutes, and $2+44=46$ so $h+m= 46\ \textbf{(A)}$.

 


	\item (\textit{AMC 12 12A 2009}) Which of the following is equal to $1 + \frac {1}{1 + \frac {1}{1 + 1}}$?


 $\textbf{(A)}\ \frac {5}{4} \qquad \textbf{(B)}\ \frac {3}{2} \qquad \textbf{(C)}\ \frac {5}{3} \qquad \textbf{(D)}\ 2 \qquad \textbf{(E)}\ 3$

 \textbf{Answer: }C



\textbf{Solution}

We compute:

 
 \begin{align*} 1 + \frac {1}{1 + \frac {1}{1 + 1}} &= 1 + \frac {1}{1 + \frac {1}{1 + 1}} \\ &= 1 + \frac {1}{1 + \frac 12} \\ &= 1 + \frac {1}{\frac 32} \\ &= 1 + \frac 23 \\ &= \frac 53 \end{align*}


 This is choice $\boxed{\text{C}}$.

 Interesting sidenote: The continued fraction $1 + \frac {1}{1 + \frac {1}{1 + 1....}}$ is equal to the golden ratio, or $\frac{1+\sqrt{5}}{2}$.

 


	\item (\textit{AMC 12 12A 2009}) What number is one third of the way from $\frac14$ to $\frac34$?


 $\textbf{(A)}\ \frac {1}{3} \qquad \textbf{(B)}\ \frac {5}{12} \qquad \textbf{(C)}\ \frac {1}{2} \qquad \textbf{(D)}\ \frac {7}{12} \qquad \textbf{(E)}\ \frac {2}{3}$

 \textbf{Answer: }B



\textbf{Solution 1}

We can rewrite the two given fractions as $\frac 3{12}$ and $\frac 9{12}$. (We multiplied all numerators and denominators by $3$.)

 Now it is obvious that the interval between them is divided into three parts by the fractions $\boxed{\frac 5{12}}$ and $\frac 7{12}$.

 



\textbf{Solution 2}

The number we seek can be obtained as a weighted average of the two endpoints, where the closer one has weight $2$ and the further one $1$. We compute:

 
 \[\dfrac{ 2\cdot\frac 14 + 1\cdot\frac 34 }3 = \dfrac{ \frac 54 }3 = \boxed{\dfrac 5{12}}\]



	\item (\textit{AMC 12 12A 2009}) Four coins are picked out of a piggy bank that contains a collection of pennies, nickels, dimes, and quarters. Which of the following could <em>not</em> be the total value of the four coins, in cents?


 $\textbf{(A)}\ 15 \qquad \textbf{(B)}\ 25 \qquad \textbf{(C)}\ 35 \qquad \textbf{(D)}\ 45 \qquad \textbf{(E)}\ 55$

 \textbf{Answer: }A



\textbf{Solution 1}

Pre-Note: This solution is kinda just guessing, idk you decide.

 We can solve this problem by trying out numbers to get the answer choices and use the process of elimination to get the answer. Thinking for a few seconds, you can see that for 15 cents there are no possible ways to get it. Now you may be inclined to chose $\textbf{A}$, but we may have missed some way to get 15 cents, so we try the other answer choices. 25 cents can be made with 3 nickels and 1 dimes, 35 cents can be made with 3 dimes and 1 nickel, 45 cents can be made with 1 quarter, 1 dime, and 2 nickels, and finally, 55 cents can be made with 1 quarter and 3 dimes. Therefore the answer is $\boxed{\textbf{A}}$

 \href{https://artofproblemsolving.com/wiki/index.php?title=User:Maxisw}{MisW} (\href{https://artofproblemsolving.com/wiki/index.php?title=User\_talk:Maxisw}{talk})Minor Edits by \href{https://artofproblemsolving.com/wiki/index.php?title=Yanchudeng&action=edit&redlink=1}{yanchudeng}

 Post-Note: This is my first time using LaTeX so it may look a little ugly.

 


	\item (\textit{AMC 12 12A 2009}) One dimension of a cube is increased by $1$, another is decreased by $1$, and the third is left unchanged. The volume of the new rectangular solid is $5$ less than that of the cube. What was the volume of the cube?


 $\textbf{(A)}\ 8 \qquad \textbf{(B)}\ 27 \qquad \textbf{(C)}\ 64 \qquad \textbf{(D)}\ 125 \qquad \textbf{(E)}\ 216$

 \textbf{Answer: }D



\textbf{Solution}

Let the original cube have edge length $a$. Then its volume is $a^3$.The new box has dimensions $a-1$, $a$, and $a+1$, hence its volume is $(a-1)a(a+1) = a^3-a$. The difference between the two volumes is $a$. As we are given that the difference is $5$, we have $a=5$, and the volume of the original cube was $5^3 = 125\Rightarrow\boxed{\text{(D)}}$.

 


	\item (\textit{AMC 12 12A 2009}) Suppose that $P = 2^m$ and $Q = 3^n$. Which of the following is equal to $12^{mn}$ for every pair of integers $(m,n)$?


 $\textbf{(A)}\ P^2Q \qquad \textbf{(B)}\ P^nQ^m \qquad \textbf{(C)}\ P^nQ^{2m} \qquad \textbf{(D)}\ P^{2m}Q^n \qquad \textbf{(E)}\ P^{2n}Q^m$

 \textbf{Answer: }E



\textbf{Solution}

We have $12^{mn} = (2\cdot 2\cdot 3)^{mn} = 2^{2mn} \cdot 3^{mn} = (2^m)^{2n} \cdot (3^n)^m = \boxed{\bold{E)} P^{2n} Q^m}$.

 


	\item (\textit{AMC 12 12A 2009}) The first three terms of an arithmetic sequence are $2x - 3$, $5x - 11$, and $3x + 1$ respectively. The $n$th term of the sequence is $2009$. What is $n$?


 $\textbf{(A)}\ 255 \qquad \textbf{(B)}\ 502 \qquad \textbf{(C)}\ 1004 \qquad \textbf{(D)}\ 1506 \qquad \textbf{(E)}\ 8037$

 \textbf{Answer: }B



\textbf{Solution}

As this is an arithmetic sequence, the difference must be constant: $(5x-11) - (2x-3) = (3x+1) - (5x-11)$. This solves to $x=4$. The first three terms then are $5$, $9$, and $13$. In general, the $n$th term is $1+4n$. Solving $1+4n=2009$, we get $n=\boxed{502}$.

 


	\item (\textit{AMC 12 12A 2009}) Four congruent rectangles are placed as shown. The area of the outer square is $4$ times that of the inner square. What is the ratio of the length of the longer side of each rectangle to the length of its shorter side?


 
                    \begin{figure}[H]
                        \centering
                        \includegraphics[width=0.5\textwidth]{Problem/AMC_12_12A_2009_8_Problem_Figure_1.png}
                    \end{figure}
                    
$\textbf{(A)}\ 3 \qquad \textbf{(B)}\ \sqrt {10} \qquad \textbf{(C)}\ 2 + \sqrt2 \qquad \textbf{(D)}\ 2\sqrt3 \qquad \textbf{(E)}\ 4$

 \textbf{Answer: }A



\textbf{Solution 1}

The area of the outer square is $4$ times that of the inner square.Therefore the side of the outer square is $\sqrt 4 = 2$ times that of the inner square.

 Then the shorter side of the rectangle is $1/4$ of the side of the outer square, and the longer side of the rectangle is $3/4$ of the side of the outer square, hence their ratio is $\boxed{3}$.

 



\textbf{Solution 2}

Let the side length of the smaller square be $1$, and let the smaller side of the rectangles be $y$. Since the larger square's area is four times larger than the smaller square's, the larger square's side length is $2$. $2$ is equivalent to $2y+1$, giving $y=1/2$. Then, the longer side of the rectangles is $3/2$. $\frac{\frac{3}{2}}{\frac{1}{2}}=\boxed{3}$.

 



\textbf{Solution 3}

Let the longer side length be $x$, and the shorter side be $a$.

 We have that $(x+a)^2=4(x-a)^2\implies x+a=2(x-a)\implies x+a=2x-2a\implies x=3a \implies \frac{x}{a}=3$

 Hence, the answer is $\boxed{A}\implies 3$




\textbf{Solution 4}

WLOG, let the shorter side of the rectangle be $1$, and the longer side be $x$ . Thus, the area of the larger square is $(1+x)^2$ . The area of the smaller square is $(x-1)^2$ . Thus, we have the equation $(1+x)^2 = 4 (x-1)^2$ . Solving for $x$, we get $x$ = $3$ .

 


	\item (\textit{AMC 12 12A 2009}) Suppose that $f(x+3)=3x^2 + 7x + 4$ and $f(x)=ax^2 + bx + c$. What is $a+b+c$?


 $\textbf{(A)}\ -1 \qquad \textbf{(B)}\ 0 \qquad \textbf{(C)}\ 1 \qquad \textbf{(D)}\ 2 \qquad \textbf{(E)}\ 3$

 \textbf{Answer: }D



\textbf{Solution 1}

As $f(x)=ax^2 + bx + c$, we have $f(1)=a\cdot 1^2 + b\cdot 1 + c = a+b+c$. 

 To compute $f(1)$, set $x=-2$ in the first formula. We get $f(1) = f(-2+3) = 3(-2)^2 + 7(-2) + 4 = 12 - 14 + 4 = \boxed{2}$.

 



\textbf{Solution 2}

Combining the two formulas, we know that $f(x+3) = a(x+3)^2 + b(x+3) + c$.

 We can rearrange the right hand side to $ax^2 + (6a+b)x + (9a+3b+c)$.

 Comparing coefficients we have $a=3$, $6a+b=7$, and $9a+3b+c = 4$. From the second equation we get $b=-11$, and then from the third we get $c=10$. Hence $a+b+c = 3-11+10 = \boxed{2}$.

 


	\item (\textit{AMC 12 12A 2009}) In quadrilateral $ABCD$, $AB = 5$, $BC = 17$, $CD = 5$, $DA = 9$, and $BD$ is an integer. What is $BD$?


 
                    \begin{figure}[H]
                        \centering
                        \includegraphics[width=0.5\textwidth]{Problem/AMC_12_12A_2009_10_Problem_Figure_1.png}
                    \end{figure}
                    
$\textbf{(A)}\ 11 \qquad \textbf{(B)}\ 12 \qquad \textbf{(C)}\ 13 \qquad \textbf{(D)}\ 14 \qquad \textbf{(E)}\ 15$

 \textbf{Answer: }C



\textbf{Solution}

By the \href{https://artofproblemsolving.com/wiki/index.php?title=Triangle\_inequality}{triangle inequality} we have $BD < DA + AB = 9 + 5 = 14$, and also $BD + CD > BC$, hence $BD > BC - CD = 17 - 5 = 12$.

 We get that $12 < BD < 14$, and as we know that $BD$ is an integer, we must have $BD=\boxed{13}$.

 


	\item (\textit{AMC 12 12A 2009}) The figures $F_1$, $F_2$, $F_3$, and $F_4$ shown are the first in a sequence of figures. For $n\ge3$, $F_n$ is constructed from $F_{n - 1}$ by surrounding it with a square and placing one more diamond on each side of the new square than $F_{n - 1}$ had on each side of its outside square. For example, figure $F_3$ has $13$ diamonds. How many diamonds are there in figure $F_{20}$?


 
                    \begin{figure}[H]
                        \centering
                        \includegraphics[width=0.5\textwidth]{Problem/AMC_12_12A_2009_11_Problem_Figure_1.png}
                    \end{figure}
                    
$\textbf{(A)}\ 401 \qquad \textbf{(B)}\ 485 \qquad \textbf{(C)}\ 585 \qquad \textbf{(D)}\ 626 \qquad \textbf{(E)}\ 761$

 \textbf{Answer: }E



\textbf{Solution 1}

Split $F_n$ into $4$ congruent triangles by its diagonals (like in the pictures in the problem). This shows that the number of diamonds it contains is equal to $4$ times the $(n-2)$th triangular number (i.e. the diamonds within the triangles or between the diagonals) and $4(n-1)+1$ (the diamonds on sides of the triangles or on the diagonals). The $n$th triangular number is $\frac{n(n+1)}{2}$. Putting this together for $F_{20}$ this gives:

 $\frac{4(18)(19)}{2}+4(19)+1=\boxed{761}$




\textbf{Solution 2}

Color the diamond layers alternately blue and red, starting from the outside. You'll get the following pattern:

 
                    \begin{figure}[H]
                        \centering
                        \includegraphics[width=0.5\textwidth]{Solution/AMC_12_12A_2009_11_Solution_Figure_1.png}
                    \end{figure}
                    

 In the figure $F_n$, the blue diamonds form a $n\times n$ square, and the red diamonds form a $(n-1)\times(n-1)$ square. Hence the total number of diamonds in $F_{20}$ is $20^2 + 19^2 = \boxed{761}$.

 



\textbf{Solution 3}

When constructing $F_n$ from $F_{n-1}$, we add $4(n-1)$ new diamonds. Let $d_n$ be the number of diamonds in $F_n$. We now know that $d_1=1$ and $\forall n>1:~ d_n=d_{n-1} + 4(n-1)$.

 Hence we get:
 \begin{align*} d_{20} & = d_{19} + 4\cdot 19 \\ & = d_{18} + 4\cdot 18 + 4\cdot 19 \\ & = \cdots \\ & = 1 + 4(1+2+\cdots+18+19) \\ & = 1 + 4\cdot\frac{19\cdot 20}2 \\ & = \boxed{761} \end{align*}





\textbf{Solution 4}

The sequence $\{ d_n\}$ goes $1, 5, 13, 25, 41,\dots$. The first finite differences go $4, 8, 12, 16, \dots$. The second finite differences go $4, 4, 4, \dots$, so we see that the second finite difference is constant. Thus, $d_n$ can be represented as a quadratic, $d_n = an^2 + bn + c$. However, we already know $d_1 = 1$, $d_2 = 3$, and $d_3 = 13$. Thus,
 \[a + b + c = d_1 = 1\]
 \[4a + 2b + c = d_2 = 3\]
 \[9a + 3b + c = d_3 = 13\]Solving this system for $a$, $b$, and $c$ gives $a = 2$, $b = -2$, $c = 1$. Finally, $d_n = 2n^2 - 2n + 1\implies d_{20} = \boxed{(E)761}$




\textbf{Solution 5}

We can use Newton's little formula for this problem. We continue with solution 4 until we get the second finite differences go $4, 4, 4, \dots$, Then we find the 20th term by applying the formula to get  1 (20-1 c 0) + 4 (20 - 1 c 1) + 4(20-1 c 2) = 1 (19 c 0 ) + 4 (19 c 1) + 4(19 c 2) = 761

 \begin{verbatim}
- timi821
\end{verbatim}




	\item (\textit{AMC 12 12A 2009}) How many positive integers less than $1000$ are $6$ times the sum of their digits?


 $\textbf{(A)}\ 0 \qquad \textbf{(B)}\ 1 \qquad \textbf{(C)}\ 2 \qquad \textbf{(D)}\ 4 \qquad \textbf{(E)}\ 12$

 \textbf{Answer: }B



\textbf{Solution 1}

The sum of the digits is at most $9+9+9=27$. Therefore the number is at most $6\cdot 27 = 162$. Out of the numbers $1$ to $162$ the one with the largest sum of digits is $99$, and the sum is $9+9=18$. Hence the sum of digits will be at most $18$.

 Also, each number with this property is divisible by $6$, therefore it is divisible by $3$, and thus also its sum of digits is divisible by $3$. Thus, the number is divisible by $18$.

 We only have six possibilities left for the sum of the digits: $3$, $6$, $9$, $12$, $15$, and $18$, but since the number is divisible by $18$, the digits can only add to $9$ or $18$. This leads to the integers $18$, $36$, $54$, $72$, $90$, and $108$ being possibilities. We can check to see that $\boxed{1}$ solution: the number $54$ is the only solution that satisfies the conditions in the problem.

 



\textbf{Solution 2}

We can write each integer between $1$ and $999$ inclusive as $\overline{abc}=100a+10b+c$ where $a,b,c\in\{0,1,\dots,9\}$ and $a+b+c>0$.The sum of digits of this number is $a+b+c$, hence we get the equation $100a+10b+c = 6(a+b+c)$. This simplifies to $94a + 4b - 5c = 0$. Clearly for $a>0$ there are no solutions, hence $a=0$ and we get the equation $4b=5c$. This obviously has only one valid solution $(b,c)=(5,4)$, hence the only solution is the number $54$!

 



\textbf{Solution 3}

The sum of the digits is at most $9+9+9=27$. Therefore the number is at most $6\cdot 27 = 162$. Since the number is $6$ times the sum of its digits, it must be divisible by $6$, therefore also by $3$, therefore the sum of its digits must be divisible by $3$. With this in mind we can conclude that the number must be divisible by $18$, not just by $6$. Since the number is divisible by $18$, it is also divisible by $9$, therefore the sum of its digits is divisible by $9$, therefore the number is divisible by $54$, which leaves us with $54$, $108$ and $162$. Only $54$ is $6$ times its digits, hence the answer is $\boxed{1}$.

 


	\item (\textit{AMC 12 12A 2009}) A ship sails $10$ miles in a straight line from $A$ to $B$, turns through an angle between $45^{\circ}$ and $60^{\circ}$, and then sails another $20$ miles to $C$. Let $AC$ be measured in miles. Which of the following intervals contains $AC^2$?

                    \begin{figure}[H]
                        \centering
                        \includegraphics[width=0.5\textwidth]{Problem/AMC_12_12A_2009_13_Problem_Figure_1.png}
                    \end{figure}
                    
$\textbf{(A)}\ [400,500] \qquad \textbf{(B)}\ [500,600] \qquad \textbf{(C)}\ [600,700] \qquad \textbf{(D)}\ [700,800] \qquad \textbf{(E)}\ [800,900]$

 \textbf{Answer: }D



\textbf{Solution}

\textbf{Answering the question} To answer the question we are asked, it is enough to compute $AC^2$ for two different angles, preferably for both extremes ($45$ and $60$ degrees). You can use the law of cosines to do so.

 Alternately, it is enough to compute $AC^2$ for one of the extreme angles. In case it falls inside one of the given intervals, we are done. In case it falls on the boundary between two options, we also have to argue whether our $AC^2$ is the minimal or the maximal possible value of $AC^2$.

 Below we show a complete solution in which we also show that all possible values of $AC^2$ do indeed lie in the given interval.

 \textbf{Complete solution} Let $C_1$ be the point the ship would reach if it turned $45^\circ$, and $C_2$ the point it would reach if it turned $60^\circ$. Obviously, $C_1$ is the furthest possible point from $A$, and $C_2$ is the closest possible point to $A$.

 Hence the interval of possible values for $AC^2$ is $[AC_2^2,AC_1^2]$.

 
                    \begin{figure}[H]
                        \centering
                        \includegraphics[width=0.5\textwidth]{Solution/AMC_12_12A_2009_13_Solution_Figure_1.png}
                    \end{figure}
                    

 We can find $AC_1^2$ and $AC_2^2$ as follows:

 Let $D_1$ and $D_2$ be the feet of the heights from $C_1$ and $C_2$ onto $AB$. The angles in the triangle $BD_1C_1$ are $45^\circ$, $45^\circ$, and $90^\circ$, hence $BD_1 = D_1C_1 = BC_1 / \sqrt 2$. Similarly, the angles in the triangle $BD_2C_2$ are $30^\circ$, $60^\circ$, and $90^\circ$, hence $BD_2 = BC_2 / 2$ and $D_2C_2 = BC_2 \sqrt 3 / 2$.

 
                    \begin{figure}[H]
                        \centering
                        \includegraphics[width=0.5\textwidth]{Solution/AMC_12_12A_2009_13_Solution_Figure_2.png}
                    \end{figure}
                    

                    \begin{figure}[H]
                        \centering
                        \includegraphics[width=0.5\textwidth]{Solution/AMC_12_12A_2009_13_Solution_Figure_3.png}
                    \end{figure}
                    

 Hence we get:

 
 \[AC_2^2 = AD_2^2 + D_2C_2^2 = 20^2 + (10\sqrt 3)^2 = 400 + 300 = 700\]

 \[AC_1^2 = AD_1^2 + D_1C_1^2 = (10 + 20/\sqrt 2)^2 + (20\sqrt 2)^2 = 100 + 400/\sqrt 2 + 200 + 200 = 500 + 200\sqrt 2 < 500 + 200\cdot 1.5 = 800\]

 Therefore for any valid $C$ the value $AC^2$ is surely in the interval $\boxed{ \textbf{(D)}[700,800] }$.

 \textbf{Alternate Solution} From the law of cosines, $500-400\cos120^\circ<AC^2<500-400\cos135^\circ\implies700<AC^2<500+200\sqrt{2}$.  This is essentially the same solution as above.  The answer is $\boxed{\textbf{(D)}}$.

 


	\item (\textit{AMC 12 12A 2009}) A triangle has vertices $(0,0)$, $(1,1)$, and $(6m,0)$, and the line $y = mx$ divides the triangle into two triangles of equal area. What is the sum of all possible values of $m$?


 $\textbf{(A)} - \!\frac {1}{3} \qquad \textbf{(B)} - \!\frac {1}{6} \qquad \textbf{(C)}\ \frac {1}{6} \qquad \textbf{(D)}\ \frac {1}{3} \qquad \textbf{(E)}\ \frac {1}{2}$

 \textbf{Answer: }B



\textbf{Solution}

Let's label the three points as $A=(0,0)$, $B=(1,1)$, and $C=(6m,0)$.

 Clearly, whenever the line $y=mx$ intersects the inside of the triangle, it will intersect the side $BC$. Let $D$ be the point of intersection. 

 The triangles $ABD$ and $ACD$ have the same height, which is the distance between the point $A$ and the line $BC$.Hence they have equal areas if and only if $D$ is the midpoint of $BC$.

 The midpoint of the segment $BC$ has coordinates $\left( \frac{6m+1}2, \frac 12 \right)$. This point lies on the line $y=mx$ if and only if $\frac 12 = m \cdot \frac{6m+1}2$. This simplifies to $6m^2 + m - 1 = 0$. Using Vieta's formulas, we find that the sum of the roots is $\boxed{\textbf{(B)} - \!\frac {1}{6}}$.

 
For illustration, below are pictures of the situation for $m=1.5$, $m=0.5$, $m=1/3$, and $m=-1/2$.

 
                    \begin{figure}[H]
                        \centering
                        \includegraphics[width=0.5\textwidth]{Solution/AMC_12_12A_2009_14_Solution_Figure_1.png}
                    \end{figure}
                    

                    \begin{figure}[H]
                        \centering
                        \includegraphics[width=0.5\textwidth]{Solution/AMC_12_12A_2009_14_Solution_Figure_2.png}
                    \end{figure}
                    

                    \begin{figure}[H]
                        \centering
                        \includegraphics[width=0.5\textwidth]{Solution/AMC_12_12A_2009_14_Solution_Figure_3.png}
                    \end{figure}
                    

                    \begin{figure}[H]
                        \centering
                        \includegraphics[width=0.5\textwidth]{Solution/AMC_12_12A_2009_14_Solution_Figure_4.png}
                    \end{figure}
                    

 Note : In contests when you get problems like these, your AOPS knowledge will not be enough. AOPs books are just the basic, always prep for more.

 



\textbf{Solution 2}

The line must pass through the triangle's centroid, since the line divides the triangle in half. The coordinates of the centroid are found by averaging those of the vertices. The slope of the line from the origin through the centroid is thus $\frac{\frac{1}{3}}{\frac{1+6m}{3}}$, which is equal to $m$.

 $\frac{\frac{1}{3}}{\frac{1+6m}{3}}=m$
$\frac{1}{3}=m\left(\frac{1+6m}{3}\right)$
$1=m(1+6m)$
$6m^2+m-1=0$

 Using Vieta's Formulas, the sum of the possible values of $m$ is $\boxed{\textbf{(B)}\; -\frac{1}{6}}$



	\item (\textit{AMC 12 12A 2009}) For what value of $n$ is $i + 2i^2 + 3i^3 + \cdots + ni^n = 48 + 49i$?


 Note: here $i = \sqrt { - 1}$.


 $\textbf{(A)}\ 24 \qquad \textbf{(B)}\ 48 \qquad \textbf{(C)}\ 49 \qquad \textbf{(D)}\ 97 \qquad \textbf{(E)}\ 98$

 \textbf{Answer: }D



\textbf{Solution 1}

We know that $i^x$ cycles every $4$ powers so we group the sum in $4$s. 
 \[i+2i^2+3i^3+4i^4=2-2i\]
 \[5i^5+6i^6+7i^7+8i^8=2-2i\]

 We can postulate that every group of $4$ is equal to $2-2i$.For 24 groups we thus, get $48-48i$ as our sum. We know the solution must lie nearThe next term is the $24*4+1=97$th term. This term is equal to $97i$ (first in a group of $4$ so $i^{97}=i$) and our sum is now $48+49i$ so $n=97\Rightarrow\boxed{\mathbf{D}}$ is our answer

 



\textbf{Solution 2}

Obviously, even powers of $i$ are real and odd powers of $i$ are imaginary. Hence the real part of the sum is $2i^2 + 4i^4 + 6i^6 + \ldots$, and the imaginary part is $i + 3i^3 + 5i^5 + \cdots$.

 Let's take a look at the real part first. We have $i^2=-1$, hence the real part simplifies to $-2+4-6+8-10+\cdots$.If there were an odd number of terms, we could pair them as follows: $-2 + (4-6) + (8-10) + \cdots$, hence the result would be negative. As we need the real part to be $48$, we must have an even number of terms. If we have an even number of terms, we can pair them as $(-2+4) + (-6+8) + \cdots$. Each parenthesis is equal to $2$, thus there are $24$ of them, and the last value used is $96$. This happens for $n=96$ and $n=97$. As $n=96$ is not present as an option, we may conclude that the answer is $\boxed{97}$.

 In a complete solution, we should now verify which of $n=96$ and $n=97$ will give us the correct imaginary part.

 We can rewrite the imaginary part as follows: $i + 3i^3 + 5i^5 + \cdots = i(1 + 3i^2 + 5i^4 + \cdots) = i(1 - 3 + 5 - \cdots)$. We need to obtain $(1 - 3 + 5 - \cdots) = 49$. Once again we can repeat the same reasoning: If the number of terms were even, the left hand side would be negative, thus the number of terms is odd. The left hand side can then be rewritten as $1 + (-3+5) + (-7+9) + \cdots$. We need $24$ parentheses, therefore the last value used is $97$. This happens when $n=97$ or $n=98$, and we are done.

 



\textbf{Solution 3 (Fast)}

Some may know the equation:

 
 \[\sum_{k=1}^{n}kr^{k-1}=\frac{1-(n+1)r^n+nr^{n+1}}{(1-r)^2}\]

 (For those curious, this comes from differentiating the equation for finite geometric sums)

 Using this equation, we have

 
 \[48+49i=i\frac{1-(n+1)i^n+ni^{n+1}}{(1-i)^2}\]
 \[=\frac{1-(n+1)i^n+ni^{n+1}}{-2}\]
 \[=-\frac{1}{2}+\frac{(n+1)i^n}{2}-\frac{ni^{n+1}}{2}\]

 Since the imaginary and the real part must be positive, we know that $i^{n+1}=-1$ or $i^{n+1}=-i$. By the same line of reason, $i^{n}=1,i$. This only works for $n\equiv 1 \mod 4$. Therefore, we have:

 
 \[\frac{-1+n}{2}+\frac{(n+1)i}{2}=48+49i\]

 Solving either the real or imaginary part gives $\boxed{\mathbf{(D) }97}$



	\item (\textit{AMC 12 12A 2009}) A circle with center $C$ is tangent to the positive $x$ and $y$-axes and externally tangent to the circle centered at $(3,0)$ with radius $1$. What is the sum of all possible radii of the circle with center $C$?


 $\textbf{(A)}\ 3 \qquad \textbf{(B)}\ 4 \qquad \textbf{(C)}\ 6 \qquad \textbf{(D)}\ 8 \qquad \textbf{(E)}\ 9$

 \textbf{Answer: }D



\textbf{Solution}

Let $r$ be the radius of our circle. For it to be tangent to the positive $x$ and $y$ axes, we must have $C=(r,r)$. For the circle to be externally tangent to the circle centered at $(3,0)$ with radius $1$, the distance between $C$ and $(3,0)$ must be exactly $r+1$.

 By the \href{https://artofproblemsolving.com/wiki/index.php?title=Pythagorean\_theorem}{Pythagorean theorem} the distance between $(r,r)$ and $(3,0)$ is $\sqrt{ (r-3)^2 + r^2 }$, hence we get the equation $(r-3)^2 + r^2 = (r+1)^2$.

 Simplifying, we obtain $r^2 - 8r + 8 = 0$. By \href{https://artofproblemsolving.com/wiki/index.php?title=Vieta\%27s\_formulas}{Vieta's formulas} the sum of the two roots of this equation is $\boxed{8}$.

 (We should actually solve for $r$ to verify that there are two distinct positive roots. In this case we get $r=4\pm 2\sqrt 2$. This is generally a good rule of thumb, but is not necessary as all of the available answers are integers, and the equation obviously doesn't factor as integers. You can also tell that there are two positive roots based on the visual interpretation.)

 
                    \begin{figure}[H]
                        \centering
                        \includegraphics[width=0.5\textwidth]{Solution/AMC_12_12A_2009_16_Solution_Figure_1.png}
                    \end{figure}
                    



	\item (\textit{AMC 12 12A 2009}) Let $a + ar_1 + ar_1^2 + ar_1^3 + \cdots$ and $a + ar_2 + ar_2^2 + ar_2^3 + \cdots$ be two different infinite geometric series of positive numbers with the same first term.  The sum of the first series is $r_1$, and the sum of the second series is $r_2$.  What is $r_1 + r_2$?


 $\textbf{(A)}\ 0\qquad \textbf{(B)}\ \frac {1}{2}\qquad \textbf{(C)}\ 1\qquad \textbf{(D)}\ \frac {1 + \sqrt {5}}{2}\qquad \textbf{(E)}\ 2$

 \textbf{Answer: }C



\textbf{Solution}

Using the formula for the sum of a geometric series we get that the sums of the given two sequences are $\frac a{1-r_1}$ and $\frac a{1-r_2}$. 

 Hence we have $\frac a{1-r_1} = r_1$ and $\frac a{1-r_2} = r_2$.This can be rewritten as $r_1(1-r_1) = r_2(1-r_2) = a$. 

 As we are given that $r_1$ and $r_2$ are distinct, these must be precisely the two roots of the equation $x^2 - x + a = 0$.

 Using \href{https://artofproblemsolving.com/wiki/index.php?title=Vieta\%27s\_formulas}{Vieta's formulas} we get that the sum of these two roots is $\boxed{1}$.

 


 



\textbf{Solution 2}

Using the previous solution we reach the equality $r_1(1-r_1) = r_2(1-r_2)$.

 Obviously, since $r_1 \neq r_2$, then $r_1 = 1 - r_2$ so $r_1 + r_2 = 1$.

 
-Vignesh Peddi

 



\textbf{Solution 3}

We basically have two infinite geometric series whose sum is equivalent to the common ratio. Let us have a geometric series: $b, br, br^2.....$. 

 The sum is: $\frac{b}{1-r} = r.$ Thus, $b = r-r^2$ and by Vieta's, the sum of the two possible values of $r$ ($r_1$ and $r_2$) is $1$.

 


 \textbf{Alternate Solution} Using the formula for the sum of a geometric series we get that the sums of the given two sequences are $\frac a{1-r_1}$ and $\frac a{1-r_2}$. 

 Hence we have $\frac a{1-r_1} = r_1$ and $\frac a{1-r_2} = r_2$.This can be rewritten as $r_1(1-r_1) = r_2(1-r_2) = a$.

 Which can be further rewritten as $r_1-r_1^2 = r_2-r_2^2$.Rearranging the equation we get $r_1-r_2 = r_1^2-r_2^2$.Expressing this as a difference of squares we get $r_1-r_2 = (r_1-r_2)(r_1+r_2)$.

 Dividing by like terms we finally get $r_1+r_2 = \boxed{1}$ as desired.

 


 
Note: It is necessary to check that $r_1-r_2\ne 0$, as you cannot divide by zero. As the problem states that the series are different, $r_1 \ne r_2$, and so there is no division by zero error.

 


	\item (\textit{AMC 12 12A 2009}) For $k > 0$, let $I_k = 10\ldots 064$, where there are $k$ zeros between the $1$ and the $6$.  Let $N(k)$ be the number of factors of $2$ in the prime factorization of $I_k$.  What is the maximum value of $N(k)$?


 $\textbf{(A)}\ 6\qquad \textbf{(B)}\ 7\qquad \textbf{(C)}\ 8\qquad \textbf{(D)}\ 9\qquad \textbf{(E)}\ 10$

 \textbf{Answer: }B



\textbf{Solution 1}

The number $I_k$ can be written as $10^{k+2} + 64 = 5^{k+2}\cdot 2^{k+2} + 2^6$.

 For $k\in\{1,2,3\}$ we have $I_k = 2^{k+2} \left( 5^{k+2} + 2^{4-k} \right)$. The first value in the parentheses is odd, the second one is even, hence their sum is odd and we have $N(k)=k+2\leq 5$.

 For $k>4$ we have $I_k=2^6 \left( 5^{k+2}\cdot 2^{k-4} + 1 \right)$. For $k>4$ the value in the parentheses is odd, hence $N(k)=6$.

 This leaves the case $k=4$. We have $I_4 = 2^6 \left( 5^6 + 1 \right)$. The value $5^6 + 1$ is obviously even. And as $5\equiv 1 \pmod 4$, we have $5^6 \equiv 1 \pmod 4$, and therefore $5^6 + 1 \equiv 2 \pmod 4$. Hence the largest power of $2$ that divides $5^6+1$ is $2^1$, and this gives us the desired maximum of the function $N$: $N(4) = \boxed{7}$.

 



\textbf{Solution 2}

Notice that $2$ is a prime factor of an integer $n$ if and only if $n$ is even. Therefore, given any sufficiently high positive integral value of $k$, dividing $I_k$ by $2^6$ yields a terminal digit of zero, and dividing by 2 again leaves us with $2^7 \cdot a = I_k$ where $a$ is an odd integer.Observe then that $\boxed{\textbf{(B)}7}$ must be the maximum value for $N(k)$ because whatever value we choose for $k$, $N(k)$ must be less than or equal to $7$.

 "Isn't this solution incomplete because we need to show that $N(k) = 7$ can be reached?"

 An example of 7 being reached is 1000064. 1000064 divided by $2^7=128$ is 7813.

 In fact, 1000064 is the ONLY $N(k)$ that satisfies $7$. All others are 6 or lower, because if there are more zeroes, to be divisible by 128 ($2^7$), the last 7 digits must be divisible by 128, but 64 isn't. Meanwhile, if there are less zeroes, we can test by division that they do not work.

 



\textbf{Solution 3}

As in the first solution, the number $I_k$ can be written as $10^{k+2} + 64 = 5^{k+2} 2^{k+2} + 2^6$. Factor $2^6$ out of the expression to get $I_k = 2^6(1+2^{k-4}5^{k+2})$.

 You can now easily see that $I_k$ is divisible by at least 6 factors of two, from $2^6$. Any other factors of two will come from the expression $(1+2^{k-4}5^{k+2})$. 

 Make the substitution: $x=k-4$.You now have $(1+2^{x}5^{x+6}) = (1+10^x 5^6)$
$5^6=15625$, so the expression becomes $(1+15625\cdot10^x)$ This is valid when $x>-4$. 

 Obviously, if $x$ is negative, the expression will be fractional and not contain any extra factors of two.

 If $x>0$, the value $15625\cdot10^x$ will end in a zero. When 1 is added to the expression, the expression will end in 1. Since numbers divisible by 2 end in 0,2,4,6, or 8, the expression is not divisible by 2 and will not contribute to the total.

 If $x=0$, the expression evaluates to $15626$. Dividing out twos, you find that this value is only divisible by one factor of 2.

 The six factors of two from earlier add to this factor of two to give $\textbf{(B)}\ 7\qquad$

 Solution written by coolak

 


 



\textbf{Solution 4}

Similar to the other solutions, notice that $I_k$ can be written as $10^{k+2}+64 \Rightarrow 2^{k+2}5^{k+2}+2^6$. Factoring out $2^6$ we see that

 $I_k = 2^6(2^{k-4}5^{k+2}+1)$

 Notice that for $k < 4$, $2^{k-4}$ will not be an integer, and will "steal" some $2$'s from the $2^6$. We don't want this to happen, since we want to maximize the exponent of $2$. We start by considering $k = 4$. Then

 $I_k = 2^6(*5^6+1) \Rightarrow 2^6(5^6+1)$
$5^6$ is an odd number; more specifically, it ends in $25$ (all powers of $5$ after $5^1$ end in $25$). Therefore the value in the parentheses will be some large number that ends in $26$. Considering the rules of divisibility, we find that $5^6+1$ is even, so it is divisible by $2$. Now our exponent of $2$ is at $7$. But the divisibility rule for $4$ is the last $2$ digits of the number must be divisible by $4$. We know the last digits: $26$, which is not divisible by $4$. Therefore $5^6 + 1$ is divisible by $2$, but not $4$. Testing more values of $k$, we find that for $k \ge 5$, the last digit becomes $1$, which means it is not even divisible by $2$. Therefore the highest possible exponent of $2$ that we can reach is $7 \Rightarrow \boxed{\text{B}}$.

 



\textbf{Solution 5 (LTE)}

Let $m=k+2$. $v_2(10^m+2^6)=6$ if $m>6$ and $v_2(10^m+2^6)=m$ if $m<6$. However, if $m=6$, then $v_2(10^6+2^6)=v_2(2^6(5^6+1))=6+v_2(5^6+1)$. By LTE, $v_2(5^6-1)=v_2(5-1)+v_2(5+1)+v_2(6)-1=2+1+1-1=3$. Since $v_2(5^6-1)=3$, $v_2(5^6+1)$ must equal $1$. So, the answer is $6+1=7 \Rightarrow \boxed{\text{B}}$




\textbf{Solution 6 (Easy)}

Note that if $k<4$, then $I_k$ is not divisible by $64=2^6$. However, if $k=5$, then $1000064$ is divisible by $64$. So, we can write $I_k$ for $k=5$ as $I_5=64(1+15625)=2^6(1+15625)$. For $k=5$, we see that $I_5=2^6(1+15625)=2^6(15626)$, which is divisible by $2^7$, because there is already a $2^6$ and the $15626$ only adds one more factor of $2$ (the last two digits of $15626$, $26$, are not divisible by $4$). For $k>5$ or $k\geq6$, there are $(k-5)$ zeroes trailing behind $15625$, giving $I_k=2^6(1+156250...0)$. $156250...0$ is clearly even if there is at least one zero. Then, adding the extra $1$ to the even $156250...0$ gives $156250...01$, which is odd. Thus, for $k>5$, $I_k=2^6(156250...01)$, which is only divisible by $2^6$. Therefore, the maximum value of $N(k)$ is $\boxed{\textbf{(B)} 7}$ when $k=5$.

 Written by ChristianZhang

 


	\item (\textit{AMC 12 12A 2009}) Andrea inscribed a circle inside a regular pentagon, circumscribed a circle around the pentagon, and calculated the area of the region between the two circles.  Bethany did the same with a regular heptagon (7 sides).  The areas of the two regions were $A$ and $B$, respectively.  Each polygon had a side length of $2$.  Which of the following is true?


 $\textbf{(A)}\ A = \frac {25}{49}B\qquad \textbf{(B)}\ A = \frac {5}{7}B\qquad \textbf{(C)}\ A = B\qquad \textbf{(D)}\ A$ $= \frac {7}{5}B\qquad \textbf{(E)}\ A = \frac {49}{25}B$

 \textbf{Answer: }C



\textbf{Solution}

In any regular polygon with side length $2$, consider the isosceles triangle formed by the center of the polygon $S$ and two consecutive vertices $X$ and $Y$. We are given that $XY=2$. Obviously $SX=SY=r$, where $r$ is the radius of the circumcircle. Let $T$ be the midpoint of $XY$. Then $XT=TY=1$, and $TS=\rho$, where $\rho$ is the radius of the incircle.

 Applying the \href{https://artofproblemsolving.com/wiki/index.php?title=Pythagorean\_theorem}{Pythagorean theorem} on the triangle $STX$, we get that $\rho^2 + 1 = r^2$.

 Then the area between the circumcircle and the incircle can be computed as $\pi r^2 - \pi \rho^2 = \pi r^2 - \pi (r^2 - 1) = \pi$.

 Hence $A=\pi$, $B=\pi$, and therefore $\boxed{\textbf{(C) }A=B}$.

 



\textbf{Solution 2 (Applying Basic Trig)}

Similar to the first solution, consider the isosceles triangle formed by each polygon.  If you drop an altitude to the side of each polygon, you get in both polygons a right triangle with base of $1$.  For both the pentagon and heptagon, the hypotenuse of these right triangles is the radii of the larger circles and the apothems (the altitude we dropped to the side of each polygon) are the radii of the smaller circles. 

 Label the apothem of the pentagon $A_{1}$ and the apothem of the heptagon $A_{2}$.  Label the hypotenuse in the pentagon $H_{1}$ and the hypotenuse in the heptagon $H_{2}$.

 Now, finding the angles in the triangles and applying trig functions to find these radii, we get the following:

 $A_{1}$ = $\cot 36^\circ$
$A_{2}$ = $\cot \frac {180}{7}^\circ$
$H_{1}$ = $\csc 36^\circ$
$H_{2}$ = $\csc \frac {180}{7}^\circ$

 Now, the areas in between the circles are:

 $\text{Pentagon circles}$ = $\pi((\csc 36^\circ)^2 - (\cot 36^\circ)^2)$
$\text{Heptagon circles}$ = $\pi((\csc \frac {180}{7}^\circ)^2 - (\cot \frac {180}{7}^\circ)^2)$

 Note the trig identity $(\cot\theta)^2 +1 = (\csc\theta)^2$, from which we can easily get that $(\csc\theta)^2 - (\cot\theta)^2 = 1$

 Thus, the area between the circles in both the heptagon and the pentagon are equivalent to $\pi$, and therefore the answer is $\boxed{\textbf{(C) }A=B}$.

 



\textbf{Solution 3}

Note that the answer choices all include powers of $5$ and $7$. This indicates that the answer must depend on the number of sides the two shapes have. 

 For simplicity, we solve the same problem with triangle and square instead of pentagon and heptagon. For the triangle setup, $A = \left(\frac{2}{\sqrt{3}}\right)^2 \cdot \pi - \left(\frac{1}{\sqrt{3}}\right)^2 \cdot \pi = \pi.$ For the square setup, $B = \left(\sqrt{2}\right)^2 \cdot \pi - \left(1\right)^2 \cdot \pi = \pi.$

 Since $A = B$ for the triangle-square case, it should be the same for the pentagon-heptagon case. Therefore, we choose $\boxed{\textbf{(C) } A=B }$.

 


	\item (\textit{AMC 12 12A 2009}) Convex quadrilateral $ABCD$ has $AB = 9$ and $CD = 12$.  Diagonals $AC$ and $BD$ intersect at $E$, $AC = 14$, and $\triangle AED$ and $\triangle BEC$ have equal areas.  What is $AE$?


 $\textbf{(A)}\ \frac {9}{2}\qquad \textbf{(B)}\ \frac {50}{11}\qquad \textbf{(C)}\ \frac {21}{4}\qquad \textbf{(D)}\ \frac {17}{3}\qquad \textbf{(E)}\ 6$

 \textbf{Answer: }E



\textbf{Solution 1}

Let $[ABC]$ denote the area of triangle $ABC$. $[AED] = [BEC]$, so $[ABD] = [AED] + [AEB] = [BEC] + [AEB] = [ABC]$. Since triangles $ABD$ and $ABC$ share a base, they also have the same height and thus $\overline{AB}||\overline{CD}$ and $\triangle{AEB}\sim\triangle{CED}$ with a ratio of $3: 4$. Hence, 
 \[AE = \frac {3}{7}\times AC = \frac {3}{7}\times 14 = 6\ \boxed{\textbf{(E)}}\].

 
                    \begin{figure}[H]
                        \centering
                        \includegraphics[width=0.5\textwidth]{Solution/AMC_12_12A_2009_20_Solution_Figure_1.png}
                    \end{figure}
                    




\textbf{Solution 2 (Trigonometry)}

Using the sine area formula AREA= absinc/2 on triangles $AED$ and $BEC$, as $\angle AED = \angle BEC$, we see that 

 
 \[(AE)(ED) = (BE)(EC)\quad \Longrightarrow\quad \frac {AE}{EC} = \frac {BE}{ED}.\]

 Since $\angle AEB = \angle DEC$, triangles $AEB$ and $DEC$ are similar. Their ratio is $\frac {AB}{CD} = \frac {3}{4}$.  Since $AE + EC = 14$, we must have $EC = 8$, so $AE = 6\ \textbf{(E)}$.

 



\textbf{Solution 3 (Fakesolve)}

The easiest way for the areas of the triangles to be equal would be if they were congruent \href{https://artofproblemsolving.com/wiki/index.php/Constraints\_Strategy}{[1]}. A way for that to work would be if $ABCD$ were simply an isosceles trapezoid! Since $AC = 14$ and $AE:EC = 3:4$ (look at the side lengths and you'll know why!), $\boxed{AE = 6}$




\textbf{Solution 4 (Easiest Way)}

Using the fact that $[AED] = [BEC]$ and the fact that $\triangle AEB \sim \triangle EDC$ (which should be trivial given the two equal triangles) we have that

 
 \[\frac{AE}{DC} = \frac{BE}{EC} = \frac{9}{12}\]

 We know that $DC=EC,$ so we have

 
 \[\frac{AE}{EC} = \frac{BE}{EC} = \frac{3}{4}\]

 Thus 

 
 \[\frac{AE}{EC} = \frac{3}{4}\]

 But $EC = 14 - AE$ so we have

 
 \[\frac{AE}{14 - AE} = \frac{3}{4}\]

 Simplifying gives $AE = \boxed{6}.$

 \textbf{Note} The two triangles that are equal in area imply that $AB$ is parallel to $DC$ which implies that $\angle{EAB} = \angle{CDE}$ and $\angle{EBA} = \angle{DCE}.$ Furthermore, since $\angle{AEB} = \angle{DEC}$ (vertical angles). By AAA similarity, $\triangle AEB \sim \triangle EDC.$



	\item (\textit{AMC 12 12A 2009}) Let $p(x) = x^3 + ax^2 + bx + c$, where $a$, $b$, and $c$ are complex numbers.  Suppose that


 $p(2009 + 9002\pi i) = p(2009) = p(9002) = 0$

 What is the number of nonreal zeros of $x^{12} + ax^8 + bx^4 + c$?


 $\textbf{(A)}\ 4\qquad \textbf{(B)}\ 6\qquad \textbf{(C)}\ 8\qquad \textbf{(D)}\ 10\qquad \textbf{(E)}\ 12$

 \textbf{Answer: }C



\textbf{Solution}

From the three zeroes, we have $p(x) = (x - (2009 + 9002\pi i))(x - 2009)(x - 9002)$.

 Then $p(x^4) = (x^4 - (2009 + 9002\pi i))(x^4 - 2009)(x^4 - 9002)=x^{12}+ax^8+bx^4+c$.

 Let's do each factor case by case:

 
            \begin{itemize}
        \item $x^4 - (2009 + 9002\pi i) = 0$: Clearly, all the fourth roots are going to be complex.
        \item $x^4 - 2009 = 0$: The real roots are $\pm \sqrt [4]{2009}$, and there are two complex roots.
        \item $x^4 - 9002 = 0$: The real roots are $\pm \sqrt [4]{9002}$, and there are two complex roots.
            \end{itemize}
            So the answer is $4 + 2 + 2 = 8\ \mathbf{(C)}$.

 \textbf{Alternative Thinking} Consider the graph of $x^4$. It is similar to a parabola, but with a wider "base". First examine $x^4-2009$ and $x^4-9002$. Clearly they are just being translated down some large amount. This will result in the $x$-axis being crossed twice, indicating $2$ real zeroes. From the \href{https://artofproblemsolving.com/wiki/index.php?title=Fundamental\_Theorem\_of\_Algebra}{Fundamental Theorem of Algebra} we know that a polynomial must have exactly as many roots as its highest degree, so we are left with $4-2$ or $2$ nonreal roots for each of the graphs. For the graph of $x^4-(2009+9002\pi i)$, it's not even possible to graph it on the Cartesian plane, so all $4$ roots will be nonreal. This is $2+2+4 = 8$ total nonreal roots $\Rightarrow \boxed{\text{C}}$.

 


	\item (\textit{AMC 12 12A 2009}) A regular octahedron has side length $1$. A plane parallel to two of its opposite faces cuts the octahedron into the two congruent solids. The polygon formed by the intersection of the plane and the octahedron has area $\frac {a\sqrt {b}}{c}$, where $a$, $b$, and $c$ are positive integers, $a$ and $c$ are relatively prime, and $b$ is not divisible by the square of any prime. What is $a + b + c$?


 $\textbf{(A)}\ 10\qquad \textbf{(B)}\ 11\qquad \textbf{(C)}\ 12\qquad \textbf{(D)}\ 13\qquad \textbf{(E)}\ 14$

 \textbf{Answer: }E



\textbf{Solution 1}

Firstly, note that the intersection of the plane must be a hexagon. Consider the net of the octahedron. Notice that the hexagon becomes a line on the net. Also, notice that, given the parallel to the faces conditions, the line must be parallel to precisely $\frac{1}{3}$ of the sides of the net. Now, notice that, through symmetry, 2 of the hexagon's vertexes lie on the midpoint of the side of the "square" in the octahedron. In the net, the condition gives you that one of the intersections of the line with the net have to be on the midpoint of the side. However, if one is on the midpoint, because of the parallel conditions, all of the vertices are on the midpoint of a side. Thus, we have a regular hexagon with a side length of the midline of an equilateral triangle with side length 1, which is $\frac{1}{2}$. Thus, the answer is$\frac {3\sqrt {3}}{8}$, and $a + b + c = 14\ \mathbf{(E)}$.

 (Can somebody clarify this and provide a diagram?)

 



\textbf{Solution 2}

The plane cuts the octahedron into two congruent solids, which allows us to consider only the cut through the top half (a square pyramid). Because the cut is parallel to one side of the pyramid and must create two congruent solids, we can see that it must take the shape of a trapezoid with right, left, and top sides being 0.5 and a bottom side of 1. This is equivalent to a unit equilateral triangle with the tip cut off at the midpoint. Accordingly, the area of the polygon is $2 \cdot \frac{1}{2} \cdot \frac{\sqrt{3}}{2} \cdot \frac{3}{4}$ or $\frac{3\sqrt{3}}{8}$, and $a + b + c = 14$.

 


	\item (\textit{AMC 12 12A 2009}) Functions $f$ and $g$ are quadratic, $g(x) = - f(100 - x)$, and the graph of $g$ contains the vertex of the graph of $f$. The four $x$-intercepts on the two graphs have $x$-coordinates $x_1$, $x_2$, $x_3$, and $x_4$, in increasing order, and $x_3 - x_2 = 150$. Then $x_4 - x_1 = m + n\sqrt p$, where $m$, $n$, and $p$ are positive integers, and $p$ is not divisible by the square of any prime. What is $m + n + p$?


 $\textbf{(A)}\ 602\qquad \textbf{(B)}\ 652\qquad \textbf{(C)}\ 702\qquad \textbf{(D)}\ 752 \qquad \textbf{(E)}\ 802$

 \textbf{Answer: }D



\textbf{Solution (Alcumus)}


                    \begin{figure}[H]
                        \centering
                        \includegraphics[width=0.5\textwidth]{Solution/AMC_12_12A_2009_23_Solution_Figure_1.png}
                    \end{figure}
                    

 The two quadratics are $180^{\circ}$ rotations of each other about $(50,0)$. Since we are only dealing with differences of roots, we can translate them to be symmetric about $(0,0)$. Now $x_3 = - x_2 = 75$ and $x_4 = - x_1$. Say our translated versions of $f$ and $g$ are $p$ and $q$, respectively, so that $p(x) = - q( - x)$. Let $x_3 = 75$ be a root of $p$ and $x_2 = - 75$ a root of $q$ by symmetry. Note that since they each contain each other's vertex, $x_1$, $x_2$, $x_3$, and $x_4$ must be roots of alternating polynomials, so $x_1$ is a root of $p$ and $x_4$ a root of $q$

 \[p(x) = a(x - 75)(x - x_1)\]

 \[q(x) = - a(x + 75)(x + x_1)\]

 The vertex of $p(x)$ is half the sum of its roots, or $\frac {75 + x_1}{2}$. We are told that the vertex of one quadratic lies on the other, so

 
 \begin{eqnarray*} p\left(\frac {75 + x_1}{2}\right) & = & a\left(\frac {75 - x_1}{2}\right)\left(\frac { - 75 + x_1}{2}\right) \\ & = & - \frac {a}{4}(x_1 - 75)^2 \\ - \frac {a}{4}(x_1 - 75)^2 & = & q\left(\frac {75 + x_1}{2}\right) \\ & = & - a\left(\frac {x_1 + 225}{2}\right)\left(\frac {3x_1 + 75}{2}\right) \\ & = & - \frac {a}{4}(x_1 + 225)(3x_1 + 75) \end{eqnarray*}

 Let $x_1 = 75u$ and divide through by $75^2$, since it will drastically simplify computations. We know $u < - 1$ and that $(u - 1)^2 = (3u + 1)(u + 3)$, or

 
 \begin{eqnarray*} 0 & = & (3u + 1)(u + 3) - (u - 1)^2 \\ & = & 3u^2 + 10u + 3 - (u^2 - 2u + 1) \\ & = & 2u^2 + 12u + 2 \\ & = & u^2 + 6u + 1 \end{eqnarray*}

 So $u = \frac { - 6\pm\sqrt {32}}{2} = - 3\pm2\sqrt2$. Since $u < - 1$, $u = - 3 - 2\sqrt2$. 

 The answer is $x_4 - x_1 = (-x_1) - x_1 = - 150u = 450 + 300\sqrt {2}$, and $450 + 300 + 2 = 752\ \mathbf{(D)}$.

 \textbf{Note} Actually it is not necessary to solve any quadratic equations, if one utilizes the two facts about the quadratic $ax^2+bx+c$ ($a>0$) that (i) the difference of the two quadratic roots equals to $\sqrt{\Delta}/a$, and (ii) that the minimum value of a quadratic equals to $-\Delta/4a$, where $\Delta=b^2-4ac$. Here is a possible adjustment to the solution:

 Without loss of generality we may "shift" $f$,$g$,$x_1,x_2,x_3,x_4$ $50$ units to the left, then the differences of $x_i$ remain the same, $x_3$ and $x_2$ are symmetrical about $0$, so $x_2=-75$, $x_3=75$. The relationship of $f$,$g$ becomes $g(x)=-f(-x)$. So we may write:

 
 \[f(x)=a(x-b)^2 + m\]
 \[g(x)=-a(x+b)^2 - m\]

 Again without loss of generality, we can assume $a>0$ and $b>0$ (Short argument is needed here instead of the lazy "wlog"). Also, the vertex of $f$ is $(b,m)$, so $m=g(b)=-4ab^2-m$, or $m=-2ab^2 <0$.

 Since $x_2,x_4$ are roots of $f$, we have the following relationship of the roots:
 \[x_2 + x_4 = 2b\]
 \[x_4 - x_2 = \sqrt{\Delta}/a = \sqrt{-4am}/a = 2\sqrt{-m/a} = 2\sqrt{2b^2} = 2\sqrt{2}b = \sqrt{2} (x_2+x_4)\]

 So $(\sqrt{2}-1)x_4 = 75(1+\sqrt{2})$, or $x_4 = 75(1+\sqrt{2})^2 = 75(3+2\sqrt{2})$.

 Therefore $x_4-x_1=2x_4=450+300\sqrt{2}$.

 



\textbf{Solution 2}


                    \begin{figure}[H]
                        \centering
                        \includegraphics[width=0.5\textwidth]{Solution/AMC_12_12A_2009_23_Solution_Figure_2.png}
                    \end{figure}
                    $g(x) = g(50-(50-x)), f(x) = f(50+(50-x))$
$=>$ the two functions are symmetric over $x=50$
$x_3-x_2 = 150 => x_2 = -25, x_3 = 125$
$g(x) = a(x-x_1)(x-x_3) = a[2(150+m)+m]m$
$f(x) = -a(x-x_2)(x_x-x_4)$
$\space$ $\space$ $\space$ $\space$ $\space$ $\space$ $\space$ $=-a(150+m)(-(150+m))$
$\space$ $\space$ $\space$ $\space$ $\space$ $\space$ $\space$ $=a(150+m)^2$
$m^2+2m(150+m) = (150+m)^2$
$3m^2+300m = m^2+300m+150^2$
$m = 75\sqrt{2}$

 Therefore $x_4-x_1 = 3(150+m) + m = 450+300\sqrt{2}$ $\space$ $450+300+2 = \boxed{752}$

 -cassphe

 


	\item (\textit{AMC 12 12A 2009}) The \textit{tower function of twos} is defined recursively as follows: $T(1) = 2$ and $T(n + 1) = 2^{T(n)}$ for $n\ge1$. Let $A = (T(2009))^{T(2009)}$ and $B = (T(2009))^A$. What is the largest integer $k$ such that


 
 \[\underbrace{\log_2\log_2\log_2\ldots\log_2B}_{k\text{ times}}\]


 is defined?


 $\textbf{(A)}\ 2009\qquad \textbf{(B)}\ 2010\qquad \textbf{(C)}\ 2011\qquad \textbf{(D)}\ 2012\qquad \textbf{(E)}\ 2013$

 \textbf{Answer: }E



\textbf{Solution (not a proof, but good enough for AMCs)}

Testing the first two (or three) positive integers instead of 2009, $k$ seems to always be 4 more. Put E and go on to tackle \#25 :)~ dolphin7

 



\textbf{Solution}

We just look at the last three logarithms for the moment, and use the fact that $\log_2 T(k) = T(k - 1)$. We wish to find:
 \begin{align*} \log_2\log_2\log_2\left({T(2009)^{\left({T(2009)}^{T(2009)}\right)}}\right) &= \log_2(T(2009)\log_2(T(2009)\log_2 T(2009)))\\ &= \log_2(T(2009)\log_2(T(2009)T(2008)))\\ &= \log_2 \{ T(2009) [ T(2008) + T(2007) ] \} \\ \end{align*}


 
Now we realize that $T(n - 1)$ is much smaller than $T(n)$. So we approximate this, remembering we have rounded down, as:

 
 \[\log_2(T(2009)) = T(2008)\]

 We have used $3$ logarithms so far. Applying $2007$ more to the left of our expression, we get $T(1) = 2$. Then we can apply the logarithm $2$ more times, until we get to $0$. So our answer is approximately $3 + 2007 + 2 = 2012$. But we rounded down, so that means that after $2012$ logarithms we get a number slightly greater than $0$, so we can apply logarithms one more time. We can be sure it is small enough so that the logarithm can only be applied $1$ more time since $2012 + 1 = 2013$ is the largest answer choice. So the answer is $\mathbf{(E)}$.

 EDIT: Why use answer choices? @above solution should explain at the end WHY they can be sure.

 \textbf{Alternative Solution} Let $L_k(x)=\log_2\log_2...\log_2(x)$ where there are $k$ $\log_2$'s. $L_k(B)$ is defined iff (if and only if) $L_{k-1}(B) > 0$ iff $L_{k-2}(B) > 1$. Note $\log_2 T(k) = T(k - 1)$, so $L_{k-2}(T(k-2))=1$. Thus, we seek the largest $k$ such that $B > T(k-2)$. Now note that

 $T(2009)^{T(2009)^{T(2009)}} > 2^{2^{T(2009)}} = T(2011)$

 so $k=2013$ satisfies the inequality. Since it is the largest choice, it is the answer.

 \textbf{Clarification} (Note: This for the people who need clear, concise reasoning in the form of mostly words, instead of a compact proof written in set theory symbols and the like. I figured if I had trouble understanding the above solutions, others would too.)

 
We begin by contemplating what $B$ actually is. Calculating the first few values of $T(n)$, we see that it quickly becomes absolutely, massively, well, MASSIVE. So actually calculating $T(2009)$, is sadly, not an option (this is a final five problem, after all). Next we consider if it is possible to express $B$ differently. We have:

 
$B = T(2009)^{T(2009)^{T(2009)}} \Rightarrow \text{calculating.... ERROR}$

 
(I didn't actually try plug it into a computer program; the point is, it's not that easy)

 So, the only thing left to do is jump right into the problem. We start by applying the $\log$s one at a time, and seeing where it goes. We have:

 
$\log_2{\left(T(2009)^{\left(T(2009)^{T(2009)}\right)}\right)}$

 
(Uhhhhhh....)

 But wait! We can use the rule $\log_b{x^y} = y*\log_b{x}$. Then the expression becomes

 
$(T(2009)^{T(2009)}*\log_2{T(2009)} \Rightarrow (T(2009)^{T(2009)}*T(2008)$

 
(because $\log_2{T(n)} = {T(n-1)}$)

 Now we try $\log$ing it again:

 
$\log_2{(T(2009)^{T(2009)}*T(2008)}$.

 
This time, we can make use of the $\log$ addition rule $\log_b{x*y} = \log_b{x}+\log_b{y}$. The expression becomes

 
$\log_2{(T(2009)^{T(2009)})} + \log_2{T(2008)}$

 
Pulling out the exponent again and further simplifying, we have

 
$\log_2{(\log_2{(B)})} = T(2009)*T(2008)+T(2007)$

 
Now we $\log$ it a third time... and we're stuck 😐. There is no pretty formula for taking the $\log$ of two things added together. So we rack our brains, sigh in frustration... and give up.

 JUST KIDDING! (Don't give up)

 Whenever you get stuck, you want to "zoom out" and re-gain perspective of what it is you're actually trying to find. The problem asks us to find the number of times we can keep taking the $\log$ of this expression, until we no longer can (once you reach a negative value, you can't take the log of it anymore; i.e. $\log$s are only defined for positive inputs). So, for the context of this problem, what we want to know is:

 
"How many times can I take the $\log$ of $B$ and get a positive value?"

 
Let's think about what $\log$ actually means. We want a number $x$, such that when $2$ is raised to the power of $x$, we get the expression in the $\log$. So, when we think about

 
$\log_2{T(2009)*T(2008)+T(2007)}$

 
What actually matters? Well, even if you have a big number, say, $124986068$, what happens when you raise $2$ to the power of that number? We would end up getting another one of these:

 
Calculating... ERROR

 
(lol 🤣)

 But in all seriousness, whenever we consider tower functions, the next tower number is always much, much bigger. So, back to the problem:

 
$\log_2{(T(2009)*T(2008)+T(2007))}$
$T(2007)$ is huge... but when you put it next to $T(2008)$, it looks like an ant. Even more so when it's next to $T(2009)*T(2008)$. So... what does that mean? Well, if we could get rid of the $T(2007)$, we could at least go another step. So... let's get rid of it! (At this point I can almost see the audience: *surprised pikachu face)

 
 \[\text{WARNING: THIS SECTION IS FOR CALCULUS USERS ONLY}\]

 (Just kidding, but yes, the following part of the explanation does incorporate some calculus-related ideas)

 
If we were to consider just

 
$\log_2{(T(2009)*T(2008))}$

 
we would get

 
$\log_2{T(2009)} + \log_2{T(2008)} \Rightarrow T(2008) + T(2007)$

 
But what if we wanted to get 

 
$T(2008) + T(2007)+1$

 
as our answer? What would the argument have to be? (In this context, "argument" refers to the value in the parentheses). Well, we can work backwards:

 
$T(2008) + T(2007)+1$
$\Rightarrow \log_2{T(2009)} + \log_2{T(2008)} + \log_2{2}$
$\Rightarrow \log_2{(T(2009)*T(2008)*2)}$

 
So the argument would have to be at least

 
$2*T(2009)*T(2008) = T(2009)*T(2008)+T(2009)*T(2008)$

 
But when we compare

 
$T(2009)*T(2008)+T(2009)*T(2008)$

 to

 $T(2009)*T(2008)+T(2007)$

 
Clearly

 
$T(2009)*T(2008) > T(2007)$

 
So,

 
$\log_2{(T(2009)*T(2008)+T(2007))}$

 
equals some value between

 
$T(2008) + T(2007)$

 and

 $T(2008) + T(2007) + 1$

 
Let's call that small excess value $\varepsilon$ (lowercase epsilon, literally "infinitely small but not 0"). We want to keep taking the $\log$ of the expression, and $\varepsilon$ is going to be pitifully small compared to $T(2008)$ or $T(2007)$. So, within the context of this problem,

 
$T(2008) + T(2007) + \varepsilon \approx T(2008) + T(2007)$

 
Now we take the $\log$ again:

 
$\log_2{(T(2008)+T(2007))}$

 
Look familiar? It should, because this is almost exactly the same scenario as the one I just made a big fuss about (Take that, you no-know-calculus noobs! On an off note, it's completely fine if you don't know calculus; AMC problems are supposed to solvable by students who don't know calculus.) So, applying the same idea as earlier we have

 
$\log_2{(T(2008)+T(2007))} \approx \log_2{T(2008)}$

 
From here on out, it's calm waters. To recap,

 
$\log_2{(\log_2{(\log_2{(B)})})} \approx T(2008)$

 
Taking the $\log$ $2007$ more times(for a total of $2010$), we get

 
$\log_2\log_2... T(2008) = T(1) = 2$

 
Taking the log again, we get

 
$\log_2{2} = 1$

 
(Now we're at $2011$ times)

 
Once more and we have

 
$\log_2{1} = 0$.

 
($2012$ times)

 
But wait, there's more! (IYKYK lol) Remember how we ignored the $\varepsilon$? Well, It's time to bring it back. Every time we took the log, we would get $T(\text{something})$ + $\varepsilon$ because we defined $\varepsilon$ to be some very, very small (but finite) number. Let's go back a few steps:

 
$\log_2{2} = 1$

 
should really be

 
$\log_2{(2+\varepsilon)} = 1+\epsilon$

 
and

 
$\log_2{1} = 0$

 
should really be

 
$\log_2{(1+\varepsilon)} = \varepsilon$

 
Since $\varepsilon$ is still a positive number, we can take the $\log$ of it one last time before it becomes negative. So, after a long and very tiring journey (for both you, the reader, and me, the writer) we have $k$ (the number of times we took the $\log$) $= 2013 \Rightarrow \boxed{\text{E}}$.~TheAsian

 


	\item (\textit{AMC 12 12A 2009}) The first two terms of a sequence are $a_1 = 1$ and $a_2 = \frac {1}{\sqrt3}$. For $n\ge1$,


 
 \[a_{n + 2} = \frac {a_n + a_{n + 1}}{1 - a_na_{n + 1}}.\]

 What is $|a_{2009}|$?


 $\textbf{(A)}\ 0\qquad \textbf{(B)}\ 2 - \sqrt3\qquad \textbf{(C)}\ \frac {1}{\sqrt3}\qquad \textbf{(D)}\ 1\qquad \textbf{(E)}\ 2 + \sqrt3$

 \textbf{Answer: }A



\textbf{Solution 1}

Consider another sequence $\{\theta_1, \theta_2, \theta_3...\}$ such that $a_n = \tan{\theta_n}$, and $0 \leq \theta_n < 180$.

 The given recurrence becomes

 
 \begin{align*} a_{n + 2} & = \frac {a_n + a_{n + 1}}{1 - a_na_{n + 1}} \\ \tan{\theta_{n + 2}} & = \frac {\tan{\theta_n} + \tan{\theta_{n + 1}}}{1 - \tan{\theta_n}\tan{\theta_{n + 1}}} \\ \tan{\theta_{n + 2}} & = \tan(\theta_{n + 1} + \theta_n) \end{align*}


 It follows that $\theta_{n + 2} \equiv \theta_{n + 1} + \theta_{n} \pmod{180}$. Since $\theta_1 = 45, \theta_2 = 30$, all terms in the sequence $\{\theta_1, \theta_2, \theta_3...\}$ will be a multiple of $15$.

 Now consider another sequence $\{b_1, b_2, b_3...\}$ such that $b_n = \theta_n/15$, and $0 \leq b_n < 12$. The sequence $b_n$ satisfies $b_1 = 3, b_2 = 2, b_{n + 2} \equiv b_{n + 1} + b_n \pmod{12}$.

 As the number of possible consecutive two terms is finite, we know that the sequence $b_n$ is periodic. Write out the first few terms of the sequence until it starts to repeat.

 $\{b_n\} = \{3,2,5,7,0,7,7,2,9,11,8,7,3,10,1,11,0,11,11,10,9,7,4,11,3,2,5,7,...\}$

 Note that $b_{25} = b_1 = 3$ and $b_{26} = b_2 = 2$. Thus $\{b_n\}$ has a period of $24$: $b_{n + 24} = b_n$.

 It follows that $b_{2009} = b_{17} = 0$ and $\theta_{2009} = 15 b_{2009} = 0$. Thus $a_{2009} = \tan{\theta_{2009}} = \tan{0} = 0.$

 Our answer is $|a_{2009}| = \boxed{\textbf{(A)}\ 0}$.

 



\textbf{Solution 2}

First, some intuition. The given recurrence relation: $a_{n + 2} = \frac {a_n + a_{n + 1}}{1 - a_na_{n + 1}}$ looks much like the tangent addition formula. So, we let $a_n=\tan{\theta_n}$. Then, we get:
 \[\tan{\theta_n}=\tan{(\theta_{n-1}+\theta_{n-2})}\]This gives us:
 \[\theta_n \equiv \theta_{n-1}+\theta_{n-2} \pmod{180}\]Now, observe that 
 \[\theta_1=45\]
 \[\theta_2=30\]We know that this sequence of values for $\theta_n$ will repeat eventually because it is mod $180$.It is just a matter of when, so we start bashing:

 
 \[\theta_3=75\]
 \[\theta_4=105\]
 \[\theta_5=0\]
 \[\theta_6=105\]
 \[\theta_7=105\]
 \[\theta_8=30\]
 \[\theta_9=135\]
 \[\theta_{10}=165\]
 \[\theta_{11}=120\]
 \[\theta_{12}=105\]
 \[\theta_{13}=45\]
 \[\theta_{14}=150\]
 \[\theta_{15}=15\]
 \[\theta_{16}=165\]
 \[\theta_{17}=0\]
 \[\theta_{18}=165\]
 \[\theta_{19}=165\]
 \[\theta_{20}=150\]
 \[\theta_{21}=135\]
 \[\theta_{22}=105\]
 \[\theta_{23}=60\]
 \[\theta_{24}=165\]
 \[\theta_{25}=45\]
 \[\theta_{26}=30\]And there is the repetition. So, this series has a period of 24. $2009 \equiv 17 \pmod{24}$, so $|a_{2009}|=|\tan{\theta_{17}}|=|\tan{0}|=|0|= \boxed{\textbf{(A)}\ 0}$ ~Firebolt360

 \textbf{Clarification} (While the above solution is quite intuitive, it's kinda hard for people who don't know many proofing symbols to understand.)

 
Notice that the formula given in the problem looks incredibly similar to the \href{https://artofproblemsolving.com/wiki/index.php?title=Trigonometric\_identities\#Angle\_Addition/Subtraction\_Identities}{Tangent Addition Formula}.Since $a_1 = 1$, let $a_1$ be $\tan{45}$. Similarly, let $a_2$ be $\tan{30}$. Then the formula for $a_3$ reads

 


 
 \[a_3 = \frac {\tan{45} + \tan{30} }{1 - \tan{45}*\tan{30}}.\]

 
But from the \href{https://artofproblemsolving.com/wiki/index.php?title=Trigonometric\_identities\#Angle\_Addition/Subtraction\_Identities}{Tangent Addition Formula} we know that this is just the formula for $\tan{(45+30)}$ or $\tan{75}$, meaning $a_3 = \tan{75}$. So, the sequence becomes a sort of \href{https://artofproblemsolving.com/wiki/index.php?title=Fibonacci\_sequence}{Fibonacci sequence}  with angle measures. We continue to sum angle measures, like so:

 
$a_4 = \tan{(75+30)} = \tan{105}$
$a_5 = \tan{(105+75)} = \tan{180}$
$a_6 = \tan{(180+105)} =$ ... wait a minute! $\tan{180} = \tan{0}$!

 
So now we have

 
$a_5 = \tan{0}$
$a_6 = \tan{(105+0)} = \tan{105}$
$a_7 = \tan{(105+105)} = \tan{210}$
$a_8 = \tan{(210+105)} =$ .... Ok this is getting tiring. Let's stop and think about whether we can simplify this. Notice that all of the angle measures are a multiple of $15$. So, let's express the angle measures as multiples of $15$.

 
Let $b_n = \frac{\arctan{(a_n)}}{15}$.

 
(Basically, $b_n$ is the angle measure of the corresponding $a_n,$ divided by $15$)

 Now we have

 
$b_1 = 3$
$b_2 = 2$
$b_3 = 5$
$b_4 = 7$
$b_5 = 12$

 
But wait... we're dealing with the $\tan$ function, which has a period (recurrence rate) of $2\pi$ or $180^{\circ}$. Since we divided the angle measures by $15$, the period is now $12$ (which aligns with what we got earlier: $a_5 = 0$). This means that we can reduce the terms of the sequence based on the $\bmod$ function, which returns the remainder after dividing by a certain amount. So now we can say $b_n \equiv b_n \bmod{12}$ ($\equiv$ denotes modular equivalence in this context). Now we continue our sequence:

 $3, 2, 5, 7, 0, 7, 7, 2, 9, 11, 8, 7, 3, 10, 1, 11, 0, 11, 11, 10, 9, 7, 4, 11, 3, 2, 5$ (It's starting to repeat now)

 So $a_1  = 3, a_{25} = 3, a_{49} = 3$, and so on. The sequence repeats every $24$ terms. The problem asks us for the value of $a_{2009}$. Let's whip out the $\bmod$ function again.

 
$2009 \bmod 24 = 17$

 
So we want the value of the $17$th number in the sequence, and looking back at it (hopefully you wrote it down somewhere) we find that $b_{17} = 0 \Rightarrow \boxed{\text{A}}$.

 


\end{enumerate}